%!TEX program = xelatex
\def\ano{2023}
\def\semestre{1}
\def\disciplina{Introdução à Álgebra Linear}
\def\turma{3}
\def\autor{Jos\'e Ant\^onio O. Freitas}
\def\instituto{MAT-UnB}

\documentclass{beamer}
\usetheme{Madrid}
\usecolortheme{beaver}
% \mode<presentation>
\usepackage[singlelinecheck=false]{caption}
\usepackage{textpos}
\usepackage{amssymb}
\usepackage{amsmath,amsfonts,amsthm,amstext}
\usepackage[justification=centering]{caption}
\usepackage[brazil]{babel}
% \usepackage[latin1]{inputenc}
\usepackage{graphicx}
\graphicspath{{../../figuras_geogebra/}{../../Pictures/}}
\usepackage{enumitem}
\usepackage{multicol}
\usepackage{answers}
\usepackage{qrcode}

\usepackage{pifont}
\newcommand{\cmark}{\ding{228}}%

% requires version 0.3 of the package
\usepackage[customcolors,beamer]{hf-tikz}

\tikzset{style green/.style={
        set fill color=green!50!lime!60,
        set border color=white,
    },
    style cyan/.style={
        set fill color=cyan!90!blue!60,
        set border color=white,
    },
    style orange/.style={
        set fill color=orange!80!red!60,
        set border color=white,
    },
    hor/.style={
        above left offset={-0.15,0.31},
        below right offset={0.15,-0.125},
        #1
    },
    ver/.style={
        above left offset={-0.1,0.3},
        below right offset={0.15,-0.15},
        #1
    }
}

\usepackage{tikz,ifthen,tikz-3dplot}
\usepackage{siunitx}
\usetikzlibrary{circuits.ee.IEC}
\usetikzlibrary{lindenmayersystems}
\usetikzlibrary[shadings]
\usepackage{pgfplots}
\pgfplotsset{compat=1.15}
\usepackage{mathrsfs}
\usetikzlibrary{arrows}
\usepackage{tcolorbox}
\newcounter{saveenumi}
\newcommand{\seti}{\setcounter{saveenumi}{\value{enumi}}}
\newcommand{\conti}{\setcounter{enumi}{\value{saveenumi}}}

\resetcounteronoverlays{saveenumi}


\renewcommand{\qedsymbol}{$\blacksquare$}

%\setbeamertemplate{itemize items}[ball]

%\setbeamertemplate{itemize items}[default]
%-------------amatrix
% Augmented matrix.  Usage (note the argument does not count the aug col):
% \begin{amatrix}{2}
    %   1  2  3 \\  4  5  6
    % \end{amatrix}
\newenvironment{amatrix}[1]{%
    \left[\begin{array}{@{}*{#1}{r}|r@{}}
    }{%
    \end{array}\right]
}

\newtheorem{definicao}{Defini\c{c}\~ao}[section]
\newtheorem{definicoes}{Defini\c{c}\~oes}[section]
\newtheorem{exemplo}{Exemplo}[section]
\newtheorem{exemplos}{Exemplos}[section]
\newtheorem{exercicio}{Exerc{\'\i}cio}
\newtheorem{observacao}{Observa{\c c}{\~a}o:}[section]
\newtheorem{observacoes}{Observa{\c c}{\~o}es:}[section]
\newtheorem*{solucao}{Solu{\c c}{\~a}o:}
\newtheorem{proposicao}{Proposi\c{c}\~ao}
\newtheorem{lema}{Lema}
\newtheorem{teorema}{Teorema}
\newtheorem{corolario}{Corol\'ario}
\newenvironment{prova}[1][Prova]{\noindent\textbf{#1:} }{\qedsymbol}%{\ \rule{0.5em}{0.5em}}
\def\ano{2023}
\def\semestre{2}
\def\turma{11}
\newcommand{\nsub}{\varsubsetneq}
\newcommand{\vaz}{\emptyset}
\newcommand{\dom}{{\rm dom\,}}
\newcommand{\im}{{\rm Im\,}}
\newcommand{\sub}{\subseteq}
\newcommand{\n}{\mathbb{N}}
\newcommand{\z}{\mathbb{Z}}
\newcommand{\rac}{\mathbb{Q}}
\newcommand{\real}{\mathbb{R}}
\newcommand{\complex}{\mathbb{C}}
\newcommand{\cp}[1]{\mathbb{#1}}
\newcommand{\ch}{\mbox{\textrm{car\,}}\nobreak}
\newcommand{\p}[1]{\mbox{\textrm{posto(#1)}}}
\newcommand{\vesp}[1]{\vspace{ #1  cm}}
\newcommand{\compcent}[1]{\vcenter{\hbox{$#1\circ$}}}
\newcommand{\comp}{\mathbin{\mathchoice
{\compcent\scriptstyle}{\compcent\scriptstyle}
{\compcent\scriptscriptstyle}{\compcent\scriptscriptstyle}}}
\newcommand\T{\rule{0pt}{2.6ex}}

\title{Transformações Lineares}
\author[\autor]{\autor}
\institute[\instituto]{\instituto}
\date{}

\begin{document}
\begin{frame}
	\maketitle
\end{frame}

\logo{\includegraphics[scale=.1]{logo-MAT.png}\vspace*{8.5cm}}

\begin{frame}
    \begin{definicao}
        Sejam $(V, +, \cdot)$ \pause e $(W, \oplus, \otimes)$ \pause espaços vetoriais sobre um corpo $\cp{K}$. \pause Uma função $T \colon V \to W$ \pause é uma \textbf{transformação linear} se:\pause
        \begin{enumerate}[label={\roman*})]
            \item $T(u_1 + u_2) \pause = T(u_1) \pause \oplus T(u_2)$ \pause para todos $u_1$, $u_2 \in V$;\pause

            \vspace*{1.5cm}

            \item $T(\lambda \cdot u) \pause = \lambda \pause \otimes T(u)$ \pause para todo $\lambda \in \cp{K}$ e todo $u \in V$.
        \end{enumerate}
    \end{definicao}
\end{frame}

\begin{frame}
    \begin{observacoes}
        \begin{enumerate}[label={\roman*})]
            \item Se $T \colon V \to V$ é uma transformação linear \pause de $V$ em $V$, então \pause $T$ é chamada de um \textbf{operador linear}.\pause

            \vspace*{1.5cm}

            \item Para simplificar a notação, \pause vamos adotar os mesmos símbolos para indicar a soma \pause e o produto por escalar \pause nos espaços vetoriais que aparecerão no decorrer do conteúdo. \pause No entanto, deve-se estar ciente que estes símbolos podem ter significados diferentes, \pause dependendo do espaço vetorial em questão.
        \end{enumerate}
    \end{observacoes}
\end{frame}

\begin{frame}
    \begin{lema}
        Sejam $V$ e $W$ espaços vetoriais sobre $\cp{K}$ \pause e $T \colon V \to W$ uma transformação linear. \pause Então:\pause
        \begin{enumerate}[label={\roman*})]\label{transformacao_linear_propriedades_basicas}
            \item $T(0_V) \pause = 0_W$, \pause onde $0_V$ e $0_W$ \pause denotam os vetores nulos de $V$ e $W$, respectivamente.\pause

            \vspace*{1cm}

            \item $T(-u) \pause = -T(u)$, \pause para cada $u \in V$.\pause

            \vspace{1cm}

            \item $T(\alpha_1u_1 \pause + \alpha_2u_2 \pause + \cdots + \pause \alpha_mu_m) \pause = \alpha_1 T(u_1) \pause + \alpha_2 T(u_2) \pause + \cdots + \alpha_m T(u_m)$, \pause onde $\alpha_i \in \cp{K}$ e $u_i \in V$ para $i = 1$, \dots, $m$.
        \end{enumerate}
    \end{lema}
\end{frame}

\begin{frame}
    \begin{lema}
        Sejam $V$ e $W$ espaços vetoriais sobre $\cp{K}$. \pause Então uma função $T \colon V \to W$ \pause é uma transformação linear se, e somente se,\pause
        \[
        T(\lambda u_1 \pause + u_2) \pause = \lambda \pause T(u_1) \pause + T(u_2),\pause
        \]
        para todos $u_1$, $u_2 \in V$ \pause e todo $\lambda \in \cp{K}$.
    \end{lema}
\end{frame}

\begin{frame}
    \begin{teorema}\label{existencia_de_transformacao_unica_dado_valores}
        Sejam $V$ e $W$ espaços vetoriais sobre $\cp{K}$. \pause Se $\{u_1, \pause \dots, \pause u_n\}$ \pause é uma base de $V$ \pause e se $\{w_1, \pause \dots, \pause w_n\} \subseteq W$, \pause então existe uma única transformação linear \pause $T \colon V \to W$ \pause tal que $T(u_1) = w_1$, \pause $T(u_2) = w_2$, \pause \dots, $T(u_n) = w_n$.
    \end{teorema}
\end{frame}

\begin{frame}
    \begin{definicao}
        Sejam $V$ e $W$ espaços vetoriais sobre um corpo $\cp{K}$ \pause e $T \colon V \to W$ uma transformação linear.\pause
        \begin{enumerate}[label={\roman*})]
            \item O conjunto\pause
            \[
            \ker T \pause = \{u \in V \pause \mid T(u) = 0_W\}\pause
            \]
            é chamado de \textbf{kernel} \pause ou \textbf{núcleo} de $T$. \pause (O núcleo de $T$ também pode ser denotado por $Nuc\ T$.)\pause

            \vspace*{1.5cm}

            \item O conjunto\pause
            \[
            \im T \pause = \{w \in W \pause \mid \mbox{ existe } \pause u \in V \pause \mbox{ tal que } \pause T(u) = w\}\pause
            \]
            é chamado de \textbf{imagem} de $T$.
        \end{enumerate}
    \end{definicao}
\end{frame}

\begin{frame}

    \begin{proposicao}
        Sejam $V$ e $W$ espaços vetoriais sobre $\cp{K}$ \pause e $T \colon V \to W$ uma transformação linear. \pause Então:\pause
        \begin{enumerate}[label={\roman*})]
            \item $\ker T$ \pause é um subespaço de $V$;\pause

            \vspace*{1cm}

            \item $\im T$ \pause é um subespaço de $W$.
        \end{enumerate}
    \end{proposicao}
\end{frame}

\begin{frame}
    \begin{teorema}[Teorema do N\'ucleo e da Imagem]\label{teorema_do_nucleo_e_da_imagem}
        Sejam $V$ e $W$ $\cp{K}$-espa\c{c}os vetoriais \pause com $\dim_\cp{K}V$ finita. \pause Seja $T \colon V \to W$ uma transforma\c{c}\~ao linear. \pause Ent\~ao
        \[
            \dim_\cp{K}V \pause = \dim_\cp{K}\ker T \pause + \dim_\cp{K}\im T.
        \]
    \end{teorema}
\end{frame}

\begin{frame}
    \begin{definicao}
        Sejam $V$ e $W$ $\cp{K}$-espa\c{c}os vetoriais \pause e $T \colon V \to W$ uma transforma\c{c}\~ao linear.\pause
        \begin{enumerate}[label={\roman*})]
            \item Dizemos que $T$ \'e \textbf{injetora} \pause se dados $u_1$, $u_2 \in V$ \pause tais que $T(u_1) = T(u_2)$, \pause ent\~ao $u_1 = u_2$. \pause De modo equivalente, \pause se $u_1$, $u_2 \in V$ s\~ao tais que $u_1 \ne u_2$, \pause ent\~ao $T(u_1) \ne T(u_2)$.\pause

            \vspace{.5cm}

            \item Dizemos que $T$ \'e \textbf{sobrejetora} \pause se $\im T = W$. \pause Em outras palavras, \pause $T$ \'e \textbf{sobrejetora} se para todo $w \in W$, \pause existe $u \in V$ tal que $T(u) = w$.\pause
        \end{enumerate}
    \end{definicao}
\end{frame}

\begin{frame}
    \begin{proposicao}\label{caracteriza_transformacao_injetora}
        Sejam $V$ e $W$ espa\c{c}os vetoriais sobre um corpo $\cp{K}$ \pause e $T \colon V \to W$ uma transforma\c{c}\~ao linear.
        \begin{enumerate}[label={\roman*})]
            \item $T$ \'e injetora \pause se, e somente se, $\ker T = \{0_V\}$.

            \vspace{.5cm}

            \item Se $\dim V = \dim W$ \pause e $T$ é injetora, \pause então $T$ transforma bases em bases, \pause isto é, se $\mathcal{B} = \{v_1, v_2, \dots, v_n\}$ é uma base de $V$, \pause então $\{T(v_1), \pause T(v_2), \pause \dots, T(v_n)\}$ é uma base de $W$.\pause

            \vspace{.5cm}

            \item Se $\dim V = \dim W$, \pause então $T$ é injetora \pause se, e somente se, \pause $T$ é sobrejetora.
        \end{enumerate}
    \end{proposicao}
\end{frame}

\begin{frame}
    \begin{observacoes}
        Seja $T \colon V \to W$ uma transformação linear, \pause onde $\dim V < \infty$ e $\dim W < \infty$.\pause
        \begin{enumerate}[label={\roman*})]
            \item Se $\dim V > \dim W$, \pause então $T$ não pode ser injetora.\pause

            \vspace{.5cm}

            \item Se $\dim V < \dim W$, \pause então $T$ não pode ser sobrejetora.
        \end{enumerate}
    \end{observacoes}
\end{frame}

\begin{frame}
    \begin{definicao}
        Se $T \colon V \to W$ é uma transformação linear \pause \textbf{injetora e sobrejetora}, \pause então dizemos que $T$ é uma transformação linear \textbf{bijetora}, \pause ou que $T$ é um \textbf{isomorfismo} \pause do espaço $V$ no espaço $W$. \pause Neste caso, dizemos que o espaço $V$ é isomorfo ao espaço $W$ \pause e denotamos isso, escrevendo $V \cong W$.\pause
    \end{definicao}

    \begin{teorema}
        Seja $V$ um espaço vetorial sobre $\cp{K}$ \pause com $\dim V = n$, \pause então $V \cong \cp{K}^n$.
    \end{teorema}
\end{frame}

\begin{frame}
    \begin{observacoes}
        Seja $V$ um espaço vetorial sobre $\cp{K}$ \pause com $\dim V = n$. \pause Então do teorema anterior temos:\pause
        \begin{enumerate}[label={\roman*})]
            \item Se $\cp{K} = \rac$, \pause então $V \cong \rac^n$.\pause

            \vspace{.5cm}

            \item Se $\cp{K} = \real$, \pause então $V \cong \real^n$.\pause

            \vspace{.5cm}
            
            \item Se $\cp{K} = \complex$, \pause então $V \cong \complex^n$.
        \end{enumerate}
    \end{observacoes}
\end{frame}

\begin{frame}
    \begin{definicao}
        Dados espaços vetoriais $V$ e $W$ sobre um corpo $\cp{K}$ \pause e
        \begin{align*}
            S \colon &V \to W \quad \mbox{e}\quad T \colon V \to W\\
                     &u \mapsto S(u)\quad \mbox{ }\quad \quad u \mapsto T(u)
        \end{align*}\pause
        definimos a \textbf{soma de $S$ com $T$} \pause como a função \pause
        \begin{align*}
            S + T \colon &V \to W\\
                         &u \mapsto (S + T)(u) = S(u) + T(u).
        \end{align*}
    \end{definicao}
\end{frame}

\begin{frame}
    \begin{proposicao}
        Sejam $V$ e $W$ espaços vetoriais sobre um corpo $\cp{K}$, \pause $S \colon V \to W$ e $T \colon V \to W$ transformações lineares. \pause Então $S + T$ também é uma transformação linear.
    \end{proposicao}
\end{frame}

\begin{frame}
    \begin{definicao}
        Dados espaços vetoriais $V$ e $W$ sobre um corpo $\cp{K}$ \pause e
        \begin{align*}
            T \colon &V \to W\\
                     &u \mapsto T(u)
        \end{align*}\pause
        uma transformação linear. \pause Para $\alpha \in \cp{K}$, \pause definimos o \textbf{produto de $T$ por um escalar $\alpha$} \pause como a função\pause 
        \begin{align*}
            \alpha T \colon &V \to W\\
                            &u \mapsto (\alpha T)(u) = \alpha T(u).
        \end{align*}
    \end{definicao}
\end{frame}

\begin{frame}
    \begin{proposicao}
        Sejam $V$ e $W$ espaços vetoriais sobre um corpo $\cp{K}$, \pause $T \colon V \to W$ uma transformação linear. \pause Para todo $\alpha \in \cp{K}$, \pause $\alpha T$ também é uma transformação linear.
    \end{proposicao}
\end{frame}

\begin{frame}
    \begin{observacao}
        Sejam $V$ e $W$ espaços vetoriais sobre um corpo $\cp{K}$. \pause Seja
        \[
            U = Hom(V, W) \pause = \{T \colon V \to W \pause \mid T \mbox{ é uma transformação linear}\}.\pause 
        \]
        Das proposições anteriores segue que $(U, +, \cdot)$ é um espaço vetorial sobre $\cp{K}$.
    \end{observacao}
\end{frame}

\begin{frame}
    \begin{definicao}
        Dadas $S \colon V \to W$ \pause e $T \colon W \to U$ transformações lineares, \pause definimos a \textbf{composição de $T$ com $S$} \pause como a aplicação:\pause 
        \begin{align*}
            T \circ S \colon &V \to U\\
                             &v \mapsto (T\circ S)(u) = T(S(v))
        \end{align*}
    \end{definicao}
\end{frame}

\begin{frame}
    \begin{proposicao}
        Se $S \colon V \to W$ e $T \colon W \to U$ são transformações lineares, \pause então a composição de $T\circ S$ também é uma transformação linear.
    \end{proposicao}
\end{frame}


\begin{frame}
    \begin{definicao}
        Dado $T \colon V \to V$ um operador linear, \pause dizemos que $T$ é \textbf{invertível} \pause se existe outra aplicação $T^{-1} \colon V \to V$ \pause tal que para qualquer $u \in V$, \pause $T^{-1}(T(u)) = u$.
    \end{definicao}
\end{frame}

\begin{frame}
    \begin{observacoes}
        \begin{enumerate}[label={\roman*})]
            \item Seja $I = I_V = Id = Id_V \colon V \to V$ \pause a transformação linear tal que $I(u) = Id(u) = u$ \pause para todo $u \in V$. \pause Então $I = Id$ é chamado de \textbf{operador identidade}.\pause

            \vspace{.5cm}

            \item Se $T \colon V \to V$ é uma transformação linear invertível, \pause então a aplicação $T^{-1} \colon V \to V$ também é uma transformação linear.\pause

            \vspace{.5cm}

            \item Se $T \colon V \to V$ é um operador invertível, \pause então $T^{-1} \colon V \to V$ é chamado de \textbf{operador inverso de $T$}.
        \end{enumerate}
    \end{observacoes}
\end{frame}

\begin{frame}
    \begin{proposicao}
        Seja $T \colon V \to V$ um operador linear. \pause Então:\pause
        \begin{enumerate}[label={\roman*})]
            \item Se $T$ é invertível \pause e $T^{-1} \colon V \to V$ é o seu inverso, \pause então
                \[
                    (T^{-1} \circ T) = \pause (T \circ T^{-1}) = \pause Id.\pause
                \]

            \vspace{.5cm}

            \item Se $V$ tem dimensão finita, \pause então $T$ é invertível \pause se, e somente se, \pause $\ker(T) = \{0_V\}$.\pause

            \vspace{.5cm}

            \item Se $T$ é invertível, \pause então dada uma base $\{v_1, v_2, \dots, v_n\}$ de $V$, \pause o conjunto $\{T(v_1), \pause T(v_2), \pause \dots, T(v_n)\}$ \pause é uma base de $V$.
        \end{enumerate}
    \end{proposicao}
\end{frame}

\begin{frame}
    \begin{proposicao}
        Se $S \colon V \to V$ e $T \colon V \to V$ \pause são operadores invertíveis, \pause então $S \circ T$ é invertível e\pause
        \[
            (S \circ T)^{-1} \pause = T^{-1} \pause \circ S^{-1}.
        \]
    \end{proposicao}
\end{frame}

\end{document}
